\section{Translation Mellin transforms and exact Stieltjes raising}
\label{stmell:section}
The reflected formulas above close periodic convolutions. Positive
translations give a complementary ordinary integral calculus for the
same generalized Stieltjes functions. The filtered positive-integer
moments below specialize the stronger logarithmic-moment theorem in
\cite{stmell:balanced}; we give a short Laplace proof on the whole
convergence strip. The zero-Mellin endpoint then yields the raising,
digamma-reconstruction and finite-operator identities developed here.
The classical inputs are the Hurwitz integral, Bernoulli values and
Lerch identity \cite[Eqs.~25.11.25, 25.11.14, 25.11.18]{intdist:DLMFHurwitz}.
These results concern the conventional Laurent coefficients
$\gamma_n(a)$, with $\gamma_0=-\psi$.

Fix positive real numbers $a_1,\ldots,a_J$ and complex weights
$c_1,\ldots,c_J$. For an integer $r\ge1$, assume
\[
 \sum_{j=1}^J c_j a_j^\ell=0\qquad(0\le\ell<r).
\]
Write $\Delta F=\sum_jc_jF(a_j)$ and
\[
 F_n(x)=\sum_jc_j\gamma_n(x+a_j),\qquad
 Z(u)=\Delta\zeta(1+u,\cdot).
\]
The Stieltjes convention is
$Z(u)=\sum_{n\ge0}(-1)^nF_n(0)u^n/n!$.
Since the weights sum to zero, $Z$ is entire.

\begin{theorem}[A cancelled translation Mellin transform]
For $n\ge0$ and $0<\operatorname{Re}v<r$, the integral
$\int_0^\infty x^{v-1}F_n(x)\,dx$ converges absolutely and equals
\begin{equation}\label{stmell:master}
 \int_0^\infty x^{v-1}F_n(x)\,dx
 =(-1)^n\left.\partial_u^n\left\{
 \frac{\Gamma(v)\Gamma(1+u-v)}{\Gamma(1+u)}
 \Delta\zeta(1+u-v,\cdot)\right\}\right|_{u=0}.
\end{equation}
At every pole of the displayed Gamma factor the product is continued
before differentiation. The right side is jointly holomorphic locally
in $u,v$ subject to $\operatorname{Re}v>0$ and
$\operatorname{Re}v<r+\operatorname{Re}u$.
\end{theorem}

\begin{proof}
Set
\[
 Q(t)=\frac{\sum_jc_je^{-a_jt}}{1-e^{-t}}.
\]
The cancellation assumptions give $Q(t)=O(t^{r-1})$ at zero;
at infinity it decreases exponentially. The standard Hurwitz integral,
summed before continuation, gives
\[
 \sum_jc_j\zeta(s,x+a_j)
 =\frac1{\Gamma(s)}\int_0^\infty t^{s-1}e^{-xt}Q(t)\,dt
 \qquad(\operatorname{Re}s>1-r).
\]
Initially this follows from the classical formula for
$\operatorname{Re}s>1$, and continuation proves the stated extension.
The integral itself is holomorphic in the larger half-plane.
Absolute Fubini, with
$\int_0^\infty x^{v-1}e^{-xt}\,dx=\Gamma(v)t^{-v}$,
then gives
\[
 \int_0^\infty x^{v-1}\sum_jc_j\zeta(s,x+a_j)\,dx
 =\frac{\Gamma(v)\Gamma(s-v)}{\Gamma(s)}
       \Delta\zeta(s-v,\cdot),
\]
whenever $\operatorname{Re}v>0$ and
$\operatorname{Re}(s-v)>1-r$. The Gamma--zeta product is represented
by the convergent $t$-integral, so all apparent poles there are
removable. Compact subdomains admit logarithmic derivatives of every
fixed order under both integrals. Put $s=1+u$ and extract the
Stieltjes coefficient to obtain the assertion.
In particular $F_n(x)=O(x^{-r}(1+\log^n x))$ as $x\to\infty$;
this also follows directly by scaling $t=y/x$ in the Laplace integral.
There is no singularity at $x=0$, because all $a_j$ are positive.
\end{proof}

\begin{corollary}[Every resonant integer Mellin value]
Let $1\le k<r$, and let $h_d(1,1/2,\ldots,1/(k-1))$ be the complete
homogeneous symmetric polynomial, defined by
\[
 \prod_{\ell=1}^{k-1}(1-u/\ell)^{-1}
 =\sum_{d\ge0}h_d(1,1/2,\ldots,1/(k-1))u^d.
\]
For an empty alphabet, $h_0=1$ and $h_d=0$ for $d>0$.
Then, for every $n\ge0$,
\begin{equation}\label{stmell:integer}
 \int_0^\infty x^{k-1}F_n(x)\,dx
 =(-1)^{n+k-1}n!\sum_{j=0}^n
 \frac{h_{n-j}(1,1/2,\ldots,1/(k-1))}{(j+1)!}
 \Delta\zeta^{[j+1]}(1-k,\cdot).
\end{equation}
\end{corollary}
\begin{proof}
The Bernoulli identity
$\zeta(1-k,a)=-B_k(a)/k$ and $k<r$ give
$\Delta\zeta(1-k,\cdot)=0$ exactly. Gamma recurrence gives
\[
 \frac{\Gamma(k)\Gamma(1+u-k)}{\Gamma(1+u)}
 =\frac{(-1)^{k-1}}u
       \prod_{\ell=1}^{k-1}(1-u/\ell)^{-1}.
\]
The simple zero of the Hurwitz combination is at least sufficient to
cancel $1/u$; it need not have precisely order one. Expanding the two
holomorphic factors after this cancellation proves the formula.
\end{proof}

For an alternating $r$th difference choose $a_j=a+jh$ and
$c_j=(-1)^j\binom rj$, $0\le j\le r$, with $a,h>0$.
Then the numerator of $Q$ is exactly $e^{-at}(1-e^{-ht})^r$.
In particular, at $r=2,k=1,n=0$, Lerch's identity gives
\begin{equation}\label{stmell:loggamma}
 \int_0^\infty
 [\gamma_0(x+a)-2\gamma_0(x+a+h)+\gamma_0(x+a+2h)]\,dx
 =\log\Gamma(a)-2\log\Gamma(a+h)+\log\Gamma(a+2h).
\end{equation}
The sign is positive: $\gamma_0=-\psi$, and the integral also follows
by the fundamental theorem of calculus and the vanishing asymptotic
second difference of $\log\Gamma$.
More generally the same second-difference integral for index $n$ is
$(-1)^n\Delta\zeta^{[n+1]}(0,\cdot)/(n+1)$.

\begin{theorem}[An ordinary integral raising the Stieltjes index]
Only the zero-sum condition on the weights is needed. Define
\[
 \mathcal TF_n=
 \int_0^1\frac{F_n(x)-F_n(0)}x\,dx
 +\int_1^\infty\frac{F_n(x)}x\,dx.
\]
These are ordinary absolutely convergent integrals. For all $n\ge0$,
\begin{equation}\label{stmell:raising}
 \boxed{\displaystyle
 \mathcal TF_n=F_{n+1}(0)+
 \sum_{k=1}^n\frac{n!}{(n-k)!}\zeta(k+1)F_{n-k}(0).}
\end{equation}
The split at $1$ fixes the normalization; a split at $\rho>0$ subtracts
$F_n(0)\log\rho$ from the displayed answer.
\end{theorem}
\begin{proof}
The Mellin integral has Laurent expansion
$F_n(0)/v+\mathcal TF_n+O(v)$ near $v=0$. This follows by subtracting
the constant on $(0,1)$; the resulting integrals are holomorphic in a
strip about $v=0$. At general $s=1+u$, the constant term at $v=0$
of the right side of the master transform is
\[
 -\Delta\zeta'(s,\cdot)-[\gamma+\psi(s)]\Delta\zeta(s,\cdot).
\]
Now use
\[
 \gamma+\psi(1+u)
 =\sum_{k\ge1}(-1)^{k+1}\zeta(k+1)u^k.
\]
Taking $(-1)^n\partial_u^n$ at zero gives the stated finite formula.
\end{proof}

For two shifts $F_n(x)=\gamma_n(a+x)-\gamma_n(b+x)$ the first identities
are $\mathcal TF_0=F_1(0)$,
$\mathcal TF_1=F_2(0)+\zeta(2)F_0(0)$, and
$\mathcal TF_2=F_3(0)+2\zeta(2)F_1(0)+2\zeta(3)F_0(0)$.

\begin{theorem}[All Stieltjes differences from a single digamma integral]
Let the weights have sum zero. Put $d_n=F_n(0)$ and
$f(x)=F_0(x)=-\sum_jc_j\psi(x+a_j)$. For $n\ge1$ let
\[
 p_n(y)=\frac{\operatorname{Im}(y+\mathrm i\pi)^n}{\pi}
 =\sum_{k=0}^{\lfloor(n-1)/2\rfloor}
 (-1)^k\binom n{2k+1}\pi^{2k}y^{n-2k-1}
 \qquad(y\in\mathbb R).
\]
Then the following two integrals converge absolutely and give exactly
\begin{align}\label{stmell:digamma-reconstruct}
 d_n={}&\int_0^1p_n(\log x)\frac{f(x)-d_0}{x}\,dx
       +\int_1^\infty p_n(\log x)\frac{f(x)}x\,dx\notag\\
 &+\mathbf1_{2\mid n}\frac{(-1)^{n/2}\pi^n}{n+1}d_0.
\end{align}
The real-coefficient polynomial on the second line of its definition
extends the formula to complex weights; the notation involving an
imaginary part is applied only to the real variable $y$.
\end{theorem}
\begin{proof}
At $u=0$ the master transform and Gamma reflection give
\[
 \int_0^\infty x^{v-1}f(x)\,dx
 =\frac{\pi}{\sin\pi v}\Delta\zeta(1-v,\cdot),
 \qquad0<\operatorname{Re}v<1.
\]
After subtracting $d_0$ on $(0,1)$, the left side is
$d_0/v+\sum_{j\ge0}M_jv^j/j!$ in a neighborhood of zero, where
\[
 M_j=\int_0^1\log^j x\frac{f(x)-d_0}{x}\,dx
       +\int_1^\infty\log^j x\frac{f(x)}x\,dx.
\]
These integrals converge: $f(x)-d_0=O(x)$ at zero and
$f(x)=O(x^{-1})$ at infinity. All fixed logarithmic weights and a
common neighborhood of $v=0$ are controlled by these same estimates.
Multiplying by $\sin(\pi v)/\pi$ yields
\[
 \sum_{n\ge0}\frac{d_nv^n}{n!}
 =\frac{\sin\pi v}{\pi}
     \left(\frac{d_0}v+\sum_{j\ge0}\frac{M_jv^j}{j!}\right).
\]
The coefficient of $v^n$ from $d_0/v$ is zero for odd $n$ and
$(-1)^{n/2}\pi^nd_0/(n+1)!$ for even $n$.
The other coefficients give the displayed polynomial $p_n$.
\end{proof}

Thus $d_1=M_0$, $d_2=2M_1-\pi^2d_0/3$,
$d_3=3M_2-\pi^2M_0$, and
$d_4=4M_3-4\pi^2M_1+\pi^4d_0/5$.
The formula uses only digamma values under one integral and applies to
arbitrary positive shifts, not just rational shifts or adjacent shifts.
There is no appeal to an infinite series of iterated integral operators.

\begin{corollary}[Finite reconstruction by ordinary translation integrals]
For $y\ge0$ define
\[
 (Tf)(y)=\int_0^1\frac{f(y+x)-f(y)}x\,dx
              +\int_1^\infty\frac{f(y+x)}x\,dx.
\]
The space of finite linear combinations of $F_0,F_1,\ldots$ is invariant under $T$.
Every finite iterate $T^jF_0$ is therefore defined by ordinary
convergent integrals. For every $n\ge0$,
\begin{equation}\label{stmell:finite-operator}
 F_n(y)=n![z^n]\frac{e^{\gamma z}}{\Gamma(1-z)}
             \sum_{j=0}^n\frac{z^j}{j!}(T^jF_0)(y).
\end{equation}
In particular $F_1=TF_0$,
$F_2=T^2F_0-\zeta(2)F_0$, and
$F_3=T^3F_0-3\zeta(2)TF_0-2\zeta(3)F_0$.
\end{corollary}
\begin{proof}
Apply the raising identity to the shifted positive nodes $a_j+y$.
It proves the invariance assertion, and hence defines every finite
iterate without any completion of the span. In the ring of formal
power series with coefficients in that span, let
$\Phi(z)=\sum_{n\ge0}F_nz^n/n!$. The raising identity states
\[
 T\Phi=\partial_z\Phi+C(z)\Phi,
 \qquad C(z)=\sum_{k\ge1}\zeta(k+1)z^k.
\]
Put $A(z)=\exp(\sum_{k\ge2}\zeta(k)z^k/k)$.
Then $\partial_z(A\Phi)=T(A\Phi)$ and $(A\Phi)(0)=F_0$.
Coefficient recursion gives
$[z^j]A\Phi=T^jF_0/j!$. Multiplication by
$A^{-1}=e^{\gamma z}/\Gamma(1-z)$ proves the finite formula.
Only formal coefficient identities and finitely many actual
integrals are asserted; no convergence of an operator exponential is
required.
\end{proof}

\paragraph{Independent one-dimensional verification route.}
Define polynomials $P_n$ by
\[
 \frac{e^{uy}}{\Gamma(1+u)}=\sum_{n\ge0}P_n(y)u^n/n!.
\]
For example $P_0=1$, $P_1=y+\gamma$,
$P_2=(y+\gamma)^2-\zeta(2)$, and
$P_3=(y+\gamma)^3-3\zeta(2)(y+\gamma)+2\zeta(3)$.
The cancelled Hurwitz integral gives
\[
 F_n(x)=(-1)^n\int_0^\infty e^{-xt}Q(t)P_n(\log t)\,dt.
\]
Hence Mellin moments can be evaluated without nested Stieltjes
quadrature, using the ordinary integral
\[
 (-1)^n\Gamma(v)\int_0^\infty t^{-v}Q(t)P_n(\log t)\,dt.
\]
The raising integral is independently evaluated by
\[
 \mathcal TF_n=(-1)^{n+1}\int_0^\infty
 Q(t)P_n(\log t)(\gamma+\log t)\,dt,
\]
using the elementary Frullani normalization
$\int_0^1(e^{-xt}-1)\,dx/x+\int_1^\infty e^{-xt}\,dx/x
=-\gamma-\log t$. These transformed quadratures are numerical
checks of the identities, not additional general proofs.

% Literature/source boundary:
% Classical analytic input: DLMF 25.11.25 (Hurwitz Laplace integral),
% DLMF 25.11.14 (Bernoulli values), DLMF 25.11.18 (Lerch identity),
% and Gamma recurrence / the Taylor expansion of digamma at 1.
% No claim of historical priority is made. A finite targeted literature
% search does not establish that these consequences are new globally.
